\documentclass[10pt,a4paper]{letter}
\usepackage[T1]{fontenc}
\usepackage[utf8]{inputenc}
\usepackage{lmodern}
\usepackage[margin=2cm,top=1.8cm]{geometry}
\usepackage[hidelinks]{hyperref}

\signature{Quang-Vinh Dang\\
British University Vietnam\\
Hung Yen, Vietnam\\
\href{mailto:vinh.dq4@buv.edu.vn}{vinh.dq4@buv.edu.vn}\\
ORCID: \href{https://orcid.org/0000-0002-3877-8024}{0000-0002-3877-8024}}


\begin{document}

\begin{letter}{The Editors\\
\textit{The Journal of Economic Inequality}}

\opening{Dear Editors,}

I am pleased to submit the manuscript \textit{``Predicted welfare is too equal: inequality and poverty measurement with machine-learning poverty maps''} for consideration as a regular article in \textit{The Journal of Economic Inequality}.

Poverty maps built with machine learning from satellite imagery, mobile phone records and other new data sources now cover most low- and middle-income countries. They are used to target transfers and, increasingly, to describe the spatial distribution of inequality. The manuscript shows that these maps understate inequality by construction. A map reports conditional-mean predictions of welfare, and welfare is a mean-preserving spread of its conditional mean. The map is therefore Lorenz-dominated by the population distribution.

The paper makes four contributions:
\begin{itemize}\setlength{\itemsep}{0pt}
\item A general lower-bound result for all Lorenz-consistent inequality indices and for distribution-sensitive poverty measures at a fixed poverty line.
\item An exact decomposition of the gap for the mean log deviation and the Theil index. It shows that the inequality of a poverty map is an ex-ante inequality of opportunity, in the sense of Ferreira and Gignoux, with the map's features as circumstances.
\item A proof that linear recalibration of predictions for shrinkage can at most recover the inequality of the conditional mean.
\item An evaluation of simulation-based corrections.
\end{itemize}

The theory is tested with a Monte Carlo study, a design-based validation that treats 1.3 million households of the 2023 American Community Survey as a census, and an application to published satellite poverty maps for 23 African countries. Plug-in maps understate the national Gini by about a fifth with household predictors and by nearly two-thirds with area-level predictors. A residual bootstrap within bins of the prediction removes most of the national-level bias. The corrections overshoot, however, when residuals are pooled across heterogeneous areas, and none of them recovers the ranking of areas by inequality when the features do not predict within-area dispersion.

I believe the manuscript fits the journal's scope on the measurement of inequality and poverty and on related statistical and data issues. It also connects directly with the journal's literature on inequality of opportunity. The results matter for researchers and agencies who report inequality from machine-learning poverty maps.

The manuscript is original, has not been published, and is not under consideration elsewhere. I have no competing interests to declare, and the study received no funding. All data are publicly available. The replication code is available on request during review and will be deposited in a public repository before publication, in line with the journal's information deposit policy.

Thank you for considering this submission.

\closing{Yours sincerely,}

\end{letter}
\end{document}
